\subsection{Spectrum of the transpose of an endomorphism}

PROPOSITION 3. --- Let E be a Banach space and let E' be its dual. Let u be an endomorphism of E.

\begin{enumerate}
\def\labelenumi{\alph{enumi})}
\item
  We have \(\operatorname{Sp}(u) = \operatorname{Sp}(^t u)\) ;
\item
  For every \(f \in \mathcal{O}(\mathrm{Sp}(u))\) , one has \(f(^t u) = {}^t f(u)\) ;
\item
  We have \(e_{\mathrm{H}}(^{t}u)=^{t}e_{\mathrm{H}}(u)\) for every subset H of \(\operatorname{Sp}(u)\) that is open and closed.
\end{enumerate}

For an endomorphism of E to be an automorphism, it is necessary and sufficient that its transpose be an automorphism of E' (EVT, IV, p.~30, cor. 5), whence assertion a).

The map \(v \mapsto {}^t v\) is a unital continuous homomorphism of the Banach algebra \(\mathcal{L}(\mathrm{E})\) into the opposite Banach algebra of \(\mathcal{L}(\mathrm{E}')\) (EVT, IV, p.~7, prop. 8). Since the algebra \(\mathcal{O}(\mathrm{Sp}({}^t u))\) is commutative, the map \(f \mapsto {}^t f(u)\) is a unital continuous homomorphism of the algebra \(\mathcal{O}(\mathrm{Sp}({}^t u))\) into the algebra \(\mathcal{L}(\mathrm{E}')\) . This homomorphism maps the germ of the identity map of \(\mathbf{C}\) to \({}^t u\) , hence coincides with the homomorphism \(f \mapsto f({}^t u)\) (th. 5 of I, p.~74). This proves \(b\) ).

Assertion \(c\) ) follows from \(b\) ), applied to the function \(f_{\mathrm{H}}\) equal to 1 in a neighborhood of H and to 0 in a neighborhood of its complement in \(\mathrm{Sp}(u)\) .

Remark. --- Let H be an open and closed subset of \(\mathrm{Sp}(u)\) , whose associated spectral decomposition is \(\mathrm{E} = \mathrm{E}_{\mathrm{H}}(u) \oplus \widetilde{\mathrm{E}}_{\mathrm{H}}(u)\) . It follows from assertion c) that if one identifies \(E'\) with \(\mathrm{E}_{\mathrm{H}}(u)' \oplus \widetilde{\mathrm{E}}_{\mathrm{H}}(u)'\) , then we have \(\mathrm{E}_{\mathrm{H}}'(^{t}u) = \mathrm{E}_{\mathrm{H}}(u)'\) and \(\widetilde{\mathrm{E}}_{\mathrm{H}}'(^{t}u) = \widetilde{\mathrm{E}}_{\mathrm{H}}(u)'\) .

